\hypertarget{appendix-6b3d696dd902}{}\label{referencestart:appendix-6b3d696dd902}
\section{Reciprocal bases and coordinate measurements}\label{appendix-6b3d696dd902-reciprocal-bases-and-coordinate-measurements}

\hypertarget{appendix-6b3d696dd902-reciprocal-bases-and-coordinate-measurements}{}

A basis reconstructs vectors; a reciprocal basis measures their coefficients. Those are different jobs when the basis is not orthonormal.

Let basis vectors \(b_1,\ldots,b_n\) form the columns of matrix \(B\). Coordinates \(c\) reconstruct a vector by

\[
v=Bc.
\]

Recovering the coordinates means solving

\[
Bc=v,
\qquad c=B^{-1}v
\]

when \(B\) is square and nonsingular. The rows of \(B^{-1}\) are the coordinate-measurement covectors. In Euclidean coordinates, their vector representatives are the columns of

\[
B^{-\mathsf T}.
\]

These reciprocal vectors \(b^i\) satisfy

\[
b^i\cdot b_j=\delta^i_j.
\]

For the slanted basis \(b_1=(1,0)\), \(b_2=(1,1)\) and vector \(v=(3,2)\), directly dotting with the basis vectors gives \((3,5)\) and is wrong as a coordinate recovery rule. The reciprocal basis is

\[
b^1=(1,-1),\qquad b^2=(0,1),
\]

so the measurements are \(b^1\cdot v=1\) and \(b^2\cdot v=2\), exactly the coefficients that reconstruct \(v=b_1+2b_2\).

A reciprocal basis is therefore a concrete bridge from vectors to covectors. It also explains why normal transformation uses an inverse transpose: normals are measurements that must continue to annihilate tangent directions after the tangent basis changes.

In numerical code, do not build a matrix inverse merely to solve one coordinate problem. Use a factorization or direct solve appropriate to the matrix structure and condition.

Full derivation and worked trace: \hyperref[01-foundations-i-ch017-recover-a-component-with-the-right-measuring-rule]{Continuous Notes/Computing Foundations - From Bits to C++ and Systems APIs\#Recover a component with the right measuring rule}.

Related: \href{../07-Vision-XR-and-Playable-Worlds-P01/main.pdf\#nameddest=appendix-46da295c6787}{Covectors linear functionals and normal transforms (Vision, XR, and Playable Worlds, Part 1, Reference notes: Covectors linear functionals and normal transforms)}, \href{../07-Vision-XR-and-Playable-Worlds-P01/main.pdf\#nameddest=appendix-ce0341f7ea0d}{Basis coordinates and dimension (Vision, XR, and Playable Worlds, Part 1, Reference notes: Basis coordinates and dimension)}, \href{../01-Bits-Programs-and-Reliable-Systems-P01/main.pdf\#nameddest=appendix-8f6a01f284a7}{Conditioning and singularity (Bits, Programs, and Reliable Systems, Part 1, Reference notes: Conditioning and singularity)}.
